% GENERATED FILE. Edit canonical lessons, then run npm run generate:books.
% canonical-source-hash: fnv1a64:92b58ebc9432d40b
% canonical-lessons: ZH-C249-yuxi, ZH-C249-fangxue, ZH-C249-duyin, ZH-W249-du-read, ZH-C249-chengji, ZH-C249-nuli

\chapter{Hitting the Books}
\label{ch:zh-hitting-the-books}
\hlchaptermodality{249}

\begin{chapteropening}
\textbf{By the end of this chapter:} \emph{I can talk about preparing a lesson, when school finishes, how a character is said, marks and working hard.}

Complete the last lesson of chapter 249: I can talk about preparing a lesson, when school finishes, how a character is said, marks and working hard.
\end{chapteropening}

\section[yùxí]{yùxí — to prepare a lesson (before class)}

\label{lesson:ZH-C249-yuxi}



\begin{quote}
Before the new one: say the Mandarin for senior high school, then the Mandarin for a pupil's parent.
\end{quote}

\subsection*{You'll want to know: yùxí}
\emph{yùxí} — \textquotedblleft{}to prepare a lesson (before class)\textquotedblright{}.

\begin{quote}
\emph{míngtiān de kè nǐ yùxí le ma} — Have you prepared tomorrow's lesson?
\end{quote}

Its characters come later; for now it is a word for the ear.

\begin{cousinweb}
\emph{yùxí} is built from \emph{yù} 'beforehand' and \emph{xí} 'to practise', both inherited from Middle Chinese, \emph{yoH} and \emph{zip} (Guangyun, in Baxter's transcription), which glosses the first, among other senses, 'before', and the second 'to learn'.
\end{cousinweb}

\begin{grammarlens}[title={What you've built}]
One more word for saying what I do.
\end{grammarlens}

\subsection*{Your turn}
\begin{itemize}
\raggedright
  \item \emph{Say it:} \emph{yùxí}
  \item \emph{Say it:} \emph{yùxí}, once more
  \item \emph{Say it:} \emph{míngtiān de kè nǐ yùxí le ma} — Have you prepared tomorrow's lesson?
  \item \emph{Recall:} say the Mandarin for senior high school, then the Mandarin for a pupil's parent, then say \emph{yùxí} again
\end{itemize}

\begin{tcolorbox}[breakable,colback=teal!4,colframe=teal!35!black,title={Before you move on}]
What does yùxí mean? (\textquotedblleft{}To prepare a lesson (before class)\textquotedblright{}.) Say it once more.
\end{tcolorbox}
\section[fàngxué]{fàngxué — to finish school for the day; (of a school) to let out}

\label{lesson:ZH-C249-fangxue}



\begin{quote}
Before the new one: say the Mandarin for a pupil's parent, then the Mandarin for to prepare a lesson.
\end{quote}

\subsection*{You'll want to know: fàngxué}
\emph{fàngxué} — \textquotedblleft{}to finish school for the day; (of a school) to let out\textquotedblright{}.

\begin{quote}
\emph{wǒmen xiàwǔ sì diǎn fàngxué} — We finish school at four in the afternoon.
\end{quote}

Its characters come later; for now it is a word for the ear.

\begin{cousinweb}
\emph{fàngxué} is built from \emph{fàng} 'to let go' and \zh{学} 'to study', both inherited from Middle Chinese, \emph{pjangH} and \emph{haewk} (Guangyun, in Baxter's transcription), which glosses the first 'to drive out' and 'to send away'. Letting pupils go home from school is not given there.
\end{cousinweb}

\begin{grammarlens}[title={What you've built}]
One more word for saying what I do.
\end{grammarlens}

\subsection*{Your turn}
\begin{itemize}
\raggedright
  \item \emph{Say it:} \emph{fàngxué}
  \item \emph{Say it:} \emph{fàngxué}, once more
  \item \emph{Say it:} \emph{wǒmen xiàwǔ sì diǎn fàngxué} — We finish school at four in the afternoon.
  \item \emph{Recall:} say the Mandarin for a pupil's parent, then the Mandarin for to prepare a lesson, then say \emph{fàngxué} again
\end{itemize}

\begin{tcolorbox}[breakable,colback=teal!4,colframe=teal!35!black,title={Before you move on}]
What does fàngxué mean? (\textquotedblleft{}To finish school for the day; (of a school) to let out\textquotedblright{}.) Say it once more.
\end{tcolorbox}
\section[\texorpdfstring{\zh{读音} (dúyīn)}{读音 (dúyīn)}]{\zh{读音} (dúyīn) — pronunciation; how a character is read}

\label{lesson:ZH-C249-duyin}



\begin{quote}
Before the new one: say the Mandarin for to prepare a lesson, then the Mandarin for to finish school for the day.
\end{quote}

\subsection*{You'll want to know: \zh{读音}}
\textbf{\zh{读音}} — \emph{dúyīn} — \textquotedblleft{}pronunciation; how a character is read\textquotedblright{}.

\begin{quote}
\textbf{\zh{这个字的读音是什么}} \emph{zhège zì de dúyīn shì shénme} — How do you say this character?
\end{quote}

The new character, in pieces: \zh{读} — ten strokes: the speech piece \zh{讠} on the left; on the right a short \zh{十} over \zh{买}. The next lesson puts it in your hand.

\begin{cousinweb}
\textbf{\zh{读音}} (\emph{dúyīn}) is built from \zh{读} 'to read aloud' and \zh{音} 'a sound', both inherited from Middle Chinese, \emph{duwk} and \emph{'im} (Guangyun, in Baxter's transcription), which glosses the first 'to recite' and quotes an older dictionary for the second: 'sound'.
\end{cousinweb}

\begin{grammarlens}[title={What you've built}]
One more word for saying what I do.
\end{grammarlens}

\subsection*{Your turn}
\begin{itemize}
\raggedright
  \item \emph{Say it:} \emph{dúyīn}
  \item \emph{Say it:} \emph{dúyīn}, once more
  \item \emph{Say it:} \emph{zhège zì de dúyīn shì shénme} — How do you say this character?
  \item \emph{Recall:} say the Mandarin for to prepare a lesson, then the Mandarin for to finish school for the day, then say \emph{dúyīn} again
\end{itemize}

\begin{tcolorbox}[breakable,colback=teal!4,colframe=teal!35!black,title={Before you move on}]
What does \zh{读音} mean? (\textquotedblleft{}Pronunciation; how a character is read\textquotedblright{}.) Say it once more.
\end{tcolorbox}
\section[\texorpdfstring{\zh{读} (dú)}{读 (dú)}]{\zh{读} — the speech piece, then \zh{十} over \zh{买}}

\label{lesson:ZH-W249-du-read}



\begin{quote}
The word first, then the character inside it.

\begin{itemize}
\raggedright
  \item \emph{Recall:} say \emph{dúyīn} — \textbf{R1}, one lesson back, before it decays
\end{itemize}

\begin{quote}
\zh{说} \zh{音}
\end{quote}

\emph{Read it:} each one aloud, with nothing said first — \emph{shuō}, \emph{yīn}
\end{quote}

\subsection*{Script — the speech piece, then \zh{十} over \zh{买}}
\hlblockfigure{\detokenize{figures/ZH-W249-du-read-filmstrip.pdf}}{How \zh{读} is written, stroke by stroke}

\begin{quote}
\zh{读}
\end{quote}

\textbf{Two pieces, left before right.}

\textbf{The left piece, \zh{讠}, two strokes,} as in \zh{说} and \zh{请}: a dot, then a short bar to the right that turns down and flicks up to the right, all one stroke.

\textbf{The right piece, eight strokes, top to bottom.} First a short \zh{十}: a bar and a vertical through it. Then \zh{买}, six strokes, as you wrote it: a bar that turns down at its right end and hooks, two short dots, then a long bar, a stroke falling to the left, and a short dot at the foot.

2 + 8 = ten.

\subsection*{Your turn}
\begin{itemize}
\raggedright
  \item \emph{Write it:} \zh{读}, all ten, with the model in front of you
  \item \emph{Cover:} write it again from memory, and count the pieces, not the strokes
\end{itemize}

\begin{tcolorbox}[breakable,colback=teal!4,colframe=teal!35!black,title={Before you move on}]
How many strokes in \zh{读}? (\textbf{Ten}.) Which piece is on the left? (\textbf{\zh{讠}, the speech piece}.) What sits under \zh{十}? (\textbf{\zh{买}}.)

Source: \href{https://raw.githubusercontent.com/chanind/hanzi-writer-data/68d10a4b21150cae5e1ebbd223eed289cf32d90c/data/\%E8\%AF\%BB.json}{Hanzi Writer Data \zh{读}, pinned stroke paths and medians}.
\end{tcolorbox}
\section[chéngjì]{chéngjì — results; marks}

\label{lesson:ZH-C249-chengji}



\begin{quote}
Before the new one: say the Mandarin for to finish school for the day, then the Mandarin for how a character is read.
\end{quote}

\subsection*{You'll want to know: chéngjì}
\emph{chéngjì} — \textquotedblleft{}results; marks\textquotedblright{}.

\begin{quote}
\emph{tā de chéngjì hěn hǎo} — His marks are very good.
\end{quote}

Its characters come later; for now it is a word for the ear.

\begin{cousinweb}
\emph{chéngjì} is built from \emph{chéng} 'to complete' and \emph{jì} 'a deed', both inherited from Middle Chinese, \emph{dzyeng} and \emph{tsek} (Guangyun, in Baxter's transcription), which glosses the first 'to finish' and the second 'merit' and 'a deed'.
\end{cousinweb}

\begin{grammarlens}[title={What you've built}]
One more word for saying what I do.
\end{grammarlens}

\subsection*{Your turn}
\begin{itemize}
\raggedright
  \item \emph{Say it:} \emph{chéngjì}
  \item \emph{Say it:} \emph{chéngjì}, once more
  \item \emph{Say it:} \emph{tā de chéngjì hěn hǎo} — His marks are very good.
  \item \emph{Recall:} say the Mandarin for to finish school for the day, then the Mandarin for how a character is read, then say \emph{chéngjì} again
\end{itemize}

\begin{tcolorbox}[breakable,colback=teal!4,colframe=teal!35!black,title={Before you move on}]
What does chéngjì mean? (\textquotedblleft{}Results; marks\textquotedblright{}.) Say it once more.
\end{tcolorbox}
\section[nǔlì]{nǔlì — hard-working; trying hard}

\label{lesson:ZH-C249-nuli}



\begin{quote}
Before the new one: say the Mandarin for how a character is read, then the Mandarin for marks at school.
\end{quote}

\subsection*{You'll want to know: nǔlì}
\emph{nǔlì} — \textquotedblleft{}hard-working; trying hard\textquotedblright{}.

\begin{quote}
\emph{tā xuéxí hěn nǔlì} — She works hard at her studies.
\end{quote}

Its characters come later; for now it is a word for the ear.

\begin{cousinweb}
\emph{nǔlì} is built from \emph{nǔ} and \emph{lì} 'strength', both inherited from Middle Chinese, \emph{nuX} and \emph{lik} (Guangyun, in Baxter's transcription), which glosses the first with this very word and the second 'sinew'.
\end{cousinweb}

\begin{grammarlens}[title={What you've built}]
One more word for saying what I do.
\end{grammarlens}

\subsection*{Your turn}
\begin{itemize}
\raggedright
  \item \emph{Say it:} \emph{nǔlì}
  \item \emph{Say it:} \emph{nǔlì}, once more
  \item \emph{Say it:} \emph{tā xuéxí hěn nǔlì} — She works hard at her studies.
  \item \emph{Recall:} say the Mandarin for how a character is read, then the Mandarin for marks at school, then say \emph{nǔlì} again
\end{itemize}

\begin{tcolorbox}[breakable,colback=teal!4,colframe=teal!35!black,title={Before you move on}]
What does nǔlì mean? (\textquotedblleft{}Hard-working; trying hard\textquotedblright{}.) Say it once more.
\end{tcolorbox}
